\chapter{Method development and results}

\section{Correlative microscopy}

In correlative microscopy, data from multiple analytical techniques are combined to gain insights that cannot be obtained using the methods individually.
However, a precise data registration to align the datasets from different instruments is a challenging but required initial step.

For complex, heterogeneous materials like cement and clinker, this can be a very valuable approach.
Techniques such as \ac{SEM}-based imaging provide high-resolution microstructural information but may lack the chemical sensitivity or mechanical data needed for a full characterisation.
Conversely, methods like \ac{LAICPMS} or \ac{NI} offer detailed chemical or mechanical properties but often with lower spatial resolution.

By combining these methods, it becomes possible to directly link specific mechanical or chemical properties to individual microstructural phases.
This requires a very good sample preparation to ensure the surface is suitable for multiple analyses.
Therefore, this section at first illustrates the usefulness of \ac{ArBIB} as surface preparation technique. 

This section demonstrates the benefits of this correlative approach through two specific examples.
\zcref{sec:EDX-NI} details the combination of \ac{NI} and \ac{EDX} to map mechanical properties onto the complex phase assemblage of a hydrated cement paste.
\zcref{sec:corr_laicpms} demonstrates how \ac{LAICPMS} can be correlated with \ac{SEM}-\ac{EDX} data to determine the distribution of trace elements within different clinker phases.
These workflows highlight how correlative microscopy provides a more holistic understanding of the structure-property relationships of cement-based materials.
